\section{Threats to Validity}

\subsection{Pre-specified versus post-hoc analyses}

Several analyses were added or changed after results were observed, and we flag them here in one place. (i)~The per-collection strict/lenient readout policy was introduced after observing a Choice-3 ECE of 43.9 on TREC-COVID; we therefore lead with the policy-free Noul readout and the strict-everywhere readout, and give both readouts for every collection (Supplementary Table~S2). (ii)~The $\pm2$ nDCG and $\pm2$ F1 equivalence margins, the TOST analyses, and the modified Knapp--Hartung sensitivity analysis are post-hoc and exploratory. (iii)~In RQ3, threshold-rule comparisons against top-3 and Adaptive-$k$ were fixed after the pilot; comparisons against top-2 and top-5 and the raw (non-Platt) thresholds were added after final results. (iv)~The 329-question MuSiQue subset is a retrospective sensitivity analysis on existing outputs, not a new confirmatory split. (v)~Added during revision, after the original results were known: the family-disjoint MuSiQue split, the downstream comparison with Qwen3-Reranker-8B/0.6B, the strict readout as the primary calibration row for every Choice-3 and Score-4 system, the extension of the Clef base control from three to seven collections, the 4{,}000-character Qwen3-Reranker-8B control, the prompt paraphrases, the SPLADE++ first stage and its pool-coverage check, the per-hop F1 breakdown, and the exact-match tables. The policies, thresholds, and paraphrases were fixed before the corresponding runs, but the decision to run each analysis was not. (vi)~Added in response to peer review, from existing outputs or local models: family-clustered tests as the primary MuSiQue inference, the completion range of the DL19 SPLADE++ difference, prediction intervals, the held-out calibration of the RAG probabilities, the rank/threshold mismatch and equal-size top-$k$ comparison, the judged-count table, the family-disjoint mechanism, exact-match, and hop analyses, the threshold comparison for every ranker, and the 4{,}000-character Qwen3-Reranker-8B runs on TREC-COVID, NFCorpus, and SciFact.

\subsection{Statistical dependence}

The three question formats of each decision model are answered in one call per pair and are not independent samples. Holm correction is applied within each system--baseline pair across collections, not across the 22 systems and formats (75 system--baseline pairs). The original MuSiQue split was by ID, and 671 of 1{,}000 final questions share at least one decomposition sub-question with the calibration or pilot splits. Our primary family-disjoint split removes this overlap, and our primary MuSiQue tests cluster questions by family, but the split was built after the original results were seen, and excluding every family that reaches the pilot leaves a split that is 88\% 2-hop, so its results speak mainly to 2-hop questions.

\subsection{Configurations rather than isolated effects}

Same-base comparisons differ in prompts and context budgets (4{,}000 characters for Jev, Clef, and Jeff; 448 tokens for the label-probability controls), so they compare configurations rather than isolate decision-model training. Qwen3-Reranker-4B/8B use 4-bit weights. Clef and Qwen3-235B cover subsets of collections, and the Clef base control covers Clef's seven. The main results use a single prompt per format; across four wordings of the question, raw DL calibration error spans up to 17 points for Qwen3-Reranker-0.6B, 13 for Qwen3-235B, and 20 for Jev Noul, while nDCG@10 spans at most 2.7 (Supplementary Material, \emph{Prompt sensitivity}). Latencies come from different hardware and serving stacks.

\subsection{Benchmarks and contamination}

Our main experiments use BM25 candidates; on SPLADE++ candidates Jev's DL19 deficit to the cross-encoders cannot be separated from incomplete judgments, so first-stage effects are confounded with pool coverage (Supplementary Material, \emph{Neural first stage}). Calibration is measured on judged candidates, and numeric-threshold transfer groups need not share identical relevance semantics. We retain the DL23 synthetic queries and DL22 propagated judgments. Jev is a closed model whose size, training data, and cutoff are unknown; Robust04, MuSiQue, NovelEval, and FutureQueryEval all predate its release, so contamination cannot be excluded, and the answer-masking perturbation cannot separate reasoning from memorization. The RAG comparison covers bge-reranker-v2-m3 and Qwen3-Reranker-8B/0.6B but not the API label-probability baselines, the Qwen comparisons were added after the final results, thresholds use target calibration labels, and the generator cannot abstain.

\section{Conclusion}

Large typed decision models are competitive pointwise passage rerankers: Jev is on par with Qwen3-Reranker-8B and ahead of bge-reranker-v2-m3 on average, with large variation across collections (on SPLADE++ candidates it trails on DL19 under standard scoring, a deficit that cannot be separated from incomplete judgments), and its gains over bge persist at equal passage budgets on News and Robust04, as does its Robust04 lead over Qwen3-Reranker-8B when both read the same 4{,}000 characters (Qwen3-Reranker-8B stays ahead on Touche and News). Same-base controls suggest that Clef's effectiveness comes largely from its underlying LLM: it does not differ significantly from its base on six of seven collections and is ahead only on Robust04, where the base reads less text, while the small Jeff model improves on its base; Jev's base model is unknown, so we cannot attribute its effectiveness either way.
Out of the box, the yes/no probabilities of Jev and Laya are better calibrated than raw cross-encoder scores (Jeff's are not better than bge's, and Jev's graded formats depend on the readout), but a cross-encoder calibrated on other collections with the same relevance threshold reaches similar error; their practical advantage is that no labelled source data are needed, and rewording the prompt moves raw calibration error by more than the gap to bge. Small open decision models remain weak deep rerankers under standard evaluation.
For RAG, Jev improves multi-hop answer quality over bge with all three generators, and also over Qwen3-Reranker-8B despite equal ranking effectiveness, and it more often ranks the answer-bearing final hop first; on a family-disjoint MuSiQue split the gains remain ($+2.8$ to $+16.9$ F1 over bge, 23 of 27 comparisons significant with family-clustered tests; over Qwen8B, 12 of 27), largest when few passages are selected. Thresholding its probabilities, however, does not improve over fixed top-$k$ selection. Calibration alone, in short, does not make context selection a simple threshold.
Future work should build family-disjoint multi-hop splits with more 3- and 4-hop questions, extend the downstream comparison to the API label-probability baselines, and match context budgets across models on all collections.
